% ===========================================================================
%  sviluppatore/03-architettura/03-rotte.tex
% ===========================================================================
\section{Le \num{246} rotte, e chi le registra}\label{sez:rotte}

La dashboard Flask ha \num{267} decoratori \cod{@route} che danno
\num{246} percorsi distinti --- la differenza sono le coppie
«\cod{GET} in una funzione, \cod{POST} in un'altra». Sono sparsi su undici
moduli, montati in un ordine che conta.

\begin{center}
\begin{tabular}{@{}L{36mm}R{14mm}L{86mm}@{}}
\toprule
\textbf{Modulo} & \textbf{Rotte} & \textbf{Come si registra} \\
\midrule
\rigaTenue
\file{app.py}             & \num{134} & \cod{@app.route} dentro
                                        \simbolo{create\_app()} \\
\file{community.py}       & \num{43}  & \cod{Blueprint("community")},
                                        senza \cod{url\_prefix} \\
\rigaTenue
\file{studio.py}          & \num{43}  & \cod{Blueprint("studio")} \\
\file{marketplace\_routes.py} & \num{12} & Agganciate \emph{al blueprint
                                        community}, non a uno proprio \\
\rigaTenue
\file{agent.py}           & \num{10}  & \cod{Blueprint("agent")} \\
\file{genetic.py}         & \num{8}   & \cod{Blueprint("genetic")} \\
\rigaTenue
\file{auth.py}            & \num{8}   & \cod{@app.route} dentro
                                        \simbolo{init\_auth(app)} \\
\file{ui\_ext\_a/b/e/f.py} & \num{9}  & Caricate automaticamente
                                        (\vediSez{ui-ext}) \\
\bottomrule
\end{tabular}
\end{center}

\begin{nota}{Due moduli senza rotte}
\file{live\_runners.py} e \file{autopubblica.py} non registrano niente: sono
moduli di servizio. Il primo è chiamato da \file{app.py} sotto
\percorsoHTTP{/api/runners/*} e \percorsoHTTP{/api/brokers/*}; il secondo parte come thread
all'avvio.
\end{nota}

\subsection{L'ordine di montaggio}\label{sez:montaggio}

\figuraTikz{montaggio-app}
  {Gli otto passi con cui \simbolo{create\_app()} assembla l'applicazione.}
  {montaggio-app}

\subsection{Il caricamento automatico delle estensioni v2}\label{sez:ui-ext}

L'ultimo passo di registrazione è diverso dagli altri: non c'è nessun elenco.

\begin{codiceInterno}{python}{Il meccanismo delle estensioni}{dashboard/app.py:4362}
    # Rotte aggiunte con il nuovo stile (v2): ogni modulo dashboard/ui_ext_*.py
    # espone register(app, data_layer, base_path). Tenerle in moduli propri
    # evita di far crescere ancora questo file.
    try:
        import importlib as _il, pkgutil as _pu, dashboard as _dpkg
        for _m in sorted(m.name for m in _pu.iter_modules(_dpkg.__path__)
                         if m.name.startswith("ui_ext_")):
            _il.import_module(f"dashboard.{_m}").register(app, data_layer, base_path)
    except Exception as _e:  # noqa: BLE001
        logging.getLogger("app").warning("v2 page extensions not loaded: %s", _e)
\end{codiceInterno}

\noindent
Cinque conseguenze, tutte rilevanti per chi estende:

\begin{enumerate}
  \item \textbf{Nessun registro da aggiornare.} Si crea
        \file{dashboard/ui\_ext\_<lettera>.py} con una funzione
        \simbolo{register(app, data\_layer, base\_path)} e al riavvio c'è.
  \item \textbf{Ordine alfabetico garantito} dal \cod{sorted}: \file{a},
        \file{b}, \file{e}, \file{f}. Le lettere \file{c} e \file{d} non
        esistono --- un buco nella serie, non un errore.
  \item \textbf{Rotte sull'app, non su un blueprint.} Gli endpoint sono
        globali (\cod{a\_overview}, \cod{b\_project\_detail}\ldots) e
        \cod{url\_for} li raggiunge senza prefisso.
  \item \textbf{Caricate dopo tutto il resto}, compreso
        \simbolo{init\_auth}: gli hook di queste estensioni girano dopo la
        guardia di autenticazione.
  \item \textbf{Fallimento silenzioso.} Un solo \cod{except} avvolge l'intero
        ciclo: se \file{ui\_ext\_a} solleva all'import, \emph{anche b, e, f non
        vengono caricati}, e tutto ciò che si vede è una riga di avviso nel
        registro.
\end{enumerate}

\begin{attenzione}{Il punto più fragile della dashboard}
Il punto~\num{5} è il modo tipico in cui metà dell'interfaccia sparisce dopo
una modifica apparentemente innocua. Il sintomo è una pagina che va in
\cod{404} o un pannello vuoto; la causa è un \cod{ImportError} in un altro
modulo. Quando qualcosa della v2 non c'è più, la prima cosa da cercare nel
registro è la riga \cod{v2 page extensions not loaded}.
\end{attenzione}

\subsection{Rotte statiche che coprono rotte dinamiche}

Werkzeug dà la precedenza ai segmenti statici, quindi
\percorsoHTTP{/api/marketplace/items} vince su \percorsoHTTP{/api/marketplace/<artifact\_id>}.
Funziona, ed è voluto. Ma significa che un artefatto il cui identificativo
fosse \cod{items}, \cod{mine}, \cod{acquisti}, \cod{diagnostica},
\cod{contestazioni}, \cod{video-check} o \cod{system-artifacts} sarebbe
\emph{irraggiungibile}. Lo stesso vale per \percorsoHTTP{/marketplace/sell} contro
\percorsoHTTP{/marketplace/<artifact\_id>}, e per \percorsoHTTP{/api/jobs/conclusi} contro
\percorsoHTTP{/api/jobs/<job\_id>}.

\begin{nota}{Perché oggi non è un problema}
Gli identificativi degli artefatti sono generati, non scelti dall'utente. Se
un giorno diventassero scegliibili --- uno \emph{slug} leggibile nell'URL, per
dire --- questa lista diventerebbe un elenco di parole riservate da vietare in
fase di pubblicazione.
\end{nota}

\subsection{Tre cose morte, da non seguire}

\begin{description}[leftmargin = 1.6em, style = nextline, itemsep = 5pt]
  \item[\rotta{POST}{/api/market/publish}]
    Ritirata. Il docstring dice \cod{"""RITIRATA."""} e la rotta risponde
    sempre \cod{403}, ma resta registrata e occupa il percorso. La
    pubblicazione vera è \rotta{POST}{/api/marketplace/publish}, in
    \file{community.py}.
  \item[L'endpoint \cod{auth\_google\_token}]
    Elencato fra quelli pubblici in \file{auth.py}, ma nessuna rotta ha
    quel nome. Riferimento pendente, innocuo.
  \item[Il prefisso \percorsoHTTP{/claim}]
    Elencato fra i percorsi che richiedono un account, ma nessuna rotta
    comincia con \percorsoHTTP{/claim}: la pagina delle contestazioni è
    \percorsoHTTP{/marketplace/rimborso/<claim\_id>}, già coperta dal prefisso
    \percorsoHTTP{/marketplace}. L'equivoco nasce dal template, che si chiama
    \file{claim.html}.
\end{description}

\subsection{Due template orfani}

In \file{dashboard/templates/} ci sono \num{35} file. Due non sono raggiunti
da nessuna rotta e da nessun \cod{include}:

\begin{itemize}
  \item \file{brokers.html} --- il suo JavaScript è stato estratto in
        \file{static/js/brokers.js} e l'interfaccia dei broker vive ora dentro
        \file{live.html}. Residuo di una ristrutturazione.
  \item \file{server\_runners.html} --- i runner non stanno più sul server, ma
        sul computer dell'utente. Funzionalità rimossa, template rimasto.
\end{itemize}

\noindent
Altri quattro non sono orfani ma non sono nemmeno pagine: \file{base.html} è
il genitore di tutte le \num{32} pagine, \file{\_robustness\_panel.html},
\file{\_v2\_profile\_security.html} e \file{\_v2\_studio\_bp.html} sono
frammenti inclusi. \file{serve\_account.html} è un caso a sé: non lo rende
nessuna rotta, lo rende la guardia di autenticazione
(\vediSez{permessi}).

\begin{attenzione}{\file{live.html} non è la pagina di Studio}
\file{studio.html} è servito da \percorsoHTTP{/studio} e da
\percorsoHTTP{/studio/dati-esterni}; \file{live.html} è servito da
\percorsoHTTP{/studio/runner-live}. Il nome del percorso dice Studio, la pagina è
Live. È facile sbagliare, e il blueprint che la registra si chiama
\cod{studio}.
\end{attenzione}
